% ===================================================================
%  اللوحة رقم ٤ — سن الرجولة والشيخوخة.
%  Traduction 2026 d'apres texto/io, controlee sur texto/fr, texto/en
%  et texto/es. Consigne : voir texto/ar/10-lawha-01.tex.
% ===================================================================

%%K t04-apar-1 apar
\VUsaut{4.0mm}
\VUtitre{15.0pt}{60}{اللوحة رقم ٤}
\VUsaut{1.6mm}
\VUtitre{11.4pt}{0}{سن الرجولة والشيخوخة.}
\VUsaut{3.2mm}

%%K t04-tit-1 sub
\VUcentre{11.4pt}{0}{\textit{المشهد الأول.}}
\VUsaut{1.6mm}
\VUtitre{11.4pt}{0}{وليمة العرس.}
\VUsaut{2.6mm}

%%K t04-tit-2 sub
\VUpk{10.2pt}{40}{الخريف.}
\VUsaut{3.0mm}

%%K t04-c1-01-1 p
1. --- كبر الشبان. وها هو أحدهم، يوحنا، يتزوج. نحن في \VUgras{وليمة العرس}
التي تجمع أسرتي العروسين حول مائدة واحدة.

%%K t04-c1-02-1 p
2. --- نحن في \VUgras{الخريف}، في شهر \VUgras{سبتمبر}. ولم تعرِّ ريح
\VUgras{أكتوبر} و\VUgras{نوفمبر} الباردة الريفَ بعدُ من ورقه الأخضر. وهواء
المساء لطيف عَطِر؛ ولذلك يُقدَّم العشاء تحت \VUgras{الشرفة} \textsuperscript{(8)}
المطلة على الحديقة.

%%K t04-c1-03-1 p
3. --- وفي البعد، على \VUgras{التل} \textsuperscript{(3)}، يقوم بيت والدَي
يوحنا، وهو \VUgras{دار ريفية} \textsuperscript{(1)} أنيقة مختفية في الخضرة؛
وكروم جميلة تعطي خمرًا ممتازة تكسو سفح الرابية. والكرّام يتعهدها ويرعاها.

%%K t04-c1-04-1 p
4. --- وفوق الكروم يمر سرب من \VUgras{الحجل} \textsuperscript{(2)} أفزعه اقتراب
\VUgras{حارس الصيد} \textsuperscript{(5)}. وهو، و\VUgras{بندقيته} تحت إبطه
و\VUgras{مخلاة الصيد} على كتفه، يفتش \VUgras{الأدغال} \textsuperscript{(4)} مع
\VUgras{كلبه} \textsuperscript{(6)}. يحرس الضيعة ويمنع الصيادين الخلسة من
إفناء الطرائد.

%%K t04-c1-05-1 p
5. --- وخارج \VUgras{الشرفة} تتسلق \VUgras{دالية} معرَّشة تتدلى منها
\VUgras{عناقيد} \textsuperscript{(7)} ثقيلة.

%%K t04-c1-06-1 p
6. --- وأزهار الخريف الجميلة التي تزين \VUgras{المائدة} \textsuperscript{(16)}
هي \VUgras{الأقحوان} \textsuperscript{(9)}؛ وقد وضع \VUgras{والد} \textsuperscript{(10)}
العروس واحدة منها في عروة سترته.

%%K t04-c1-07-1 p
7. --- والطعام يشارف نهايته. وقد بدأ \VUgras{الخادم} \textsuperscript{(11)} يأتي
بأطباق الحلوى، وسيرفع بعد قليل ما لم يعد له لزوم من أجزاء المائدة:
\VUgras{الملاعق} \textsuperscript{(14)} و\VUgras{الشوك} \textsuperscript{(12)}
و\VUgras{السكاكين} \textsuperscript{(13)} و\VUgras{الصحاف} \textsuperscript{(15)}.

%%K t04-tit-3 sub
\VUpk{10.2pt}{40}{الأسرة.}
\VUsaut{3.0mm}

%%K t04-c1-08-1 p
8. --- على الوجوه كلها بِشْر، والابتسامة التي تنظر بها \VUgras{أم}
\textsuperscript{(17)} العريس إلى أولادها تدل على مبلغ سعادتها.

%%K t04-c1-09-1 p
9. --- وهذا الزواج يجمع أسرتين موسِرتين من أسر الناحية. فيوحنا،
\VUgras{الزوج} \textsuperscript{(18)}، هو \VUgras{ابن} مالك أرض كبير، وقد صار
\VUgras{صهر} صاحب مصنع ناجح.

%%K t04-c1-10-1 p
10. --- و\VUgras{الزوجان} شابان، ومع ذلك فقد ظلا \VUgras{مخطوبين} زمنًا
طويلًا. و\VUgras{الزوجة الشابة} \textsuperscript{(19)}، مارتا، فاتنة في
\VUgras{ثوبها الأبيض} المزيّن بزهر النارنج.

%%K t04-c1-11-1 p
11. --- و\VUgras{حَمُوها} \textsuperscript{(20)} مغتبط بـ\VUgras{كَنّة} بهذه
الطيبة، و\VUgras{حماة} \textsuperscript{(21)} يوحنا تبتسم إذ ترى \VUgras{ابنتها}
بهذا الجمال.

%%K t04-c1-12-1 p
12. --- وفي طرف المائدة يجلس سائر \VUgras{المدعوين} \textsuperscript{(22)}:
\VUgras{أقارب وأصدقاء وأبناء عمومة}. و\VUgras{رئيس الخدم} \textsuperscript{(23)}،
ومنديله على ذراعه، يخدمهم في خفة.

%%K t04-tit-4 sub
\VUpk{10.2pt}{40}{الأصدقاء.}
\VUsaut{3.0mm}

%%K t04-c1-13-1 p
13. --- على الجانب الأيمن من المائدة ها هي \VUgras{آنسة} \textsuperscript{(24)}،
\VUgras{صديقة} للعروس، ثم \VUgras{أخو} العروس وهو \VUgras{إشبينها}
\textsuperscript{(25)}. وهو يتحدث مع \VUgras{إشبينته} \textsuperscript{(26)}، أخت
العريس، التي تصير بهذا الزواج \VUgras{سِلفة} مارتا. وهي شابة فاتنة. لها
\VUgras{حُلي} جميلة و\VUgras{عقد لؤلؤ} و\VUgras{سوار} من ذهب.

%%K t04-c1-14-1 p
14. --- وقربهم، السيد الواقف الرافع كأسه \VUgras{صديق} \textsuperscript{(27)}
للأسرة. وهو أحد \VUgras{شاهدَي العريس}. يقيم \VUgras{نخبًا} ويشرب في
\VUgras{صحة} العروسين ويتمنى لهما سعادة طويلة. وهو في ثياب رسمية: سترة
سهرة و\VUgras{حذاء لامع} \textsuperscript{(28)}. وهو يرافق \VUgras{الجدة}
\textsuperscript{(29)}، وهي \VUgras{أرملة}.

%%K t04-c1-15-1 p
15. --- والشاب ذو \VUgras{سترة السهرة} \textsuperscript{(31)} والشارب الأسود
الرفيع هو \VUgras{صهر} \textsuperscript{(30)} يوحنا. وهو مهندس مرموق. يجلس إلى
جانب \VUgras{سيدة} \textsuperscript{(33)} أنيقة جدًا، هي \VUgras{أخت مارتا الكبرى}
وأم تلك الطفلة اللطيفة المقبلة على ذراع ابن عمها لتقدم زهرة
و\VUgras{تتمنى} \VUgras{مساء الخير}. ولهذه السيدة \VUgras{قرطان} جميلان
جدًا و\VUgras{تصفيفة شعر} أنيقة.

%%K t04-c1-16-1 p
16. --- وابن العم الصغير، وما أغربه في \VUgras{مِرْيَلته} \textsuperscript{(36)}
و\VUgras{سرواله} \textsuperscript{(37)} الواسع، جدير بالشفقة. إنه \VUgras{يتيم}
\textsuperscript{(35)}. فقد أباه وأمه وهو صغير جدًا. وعمه هو \VUgras{وصيّه}
\textsuperscript{(38)}، وقد بقي عزبًا ليربي الطفل المسكين. وهو رجل طيب. أصلع
بعض الشيء وقصير النظر جدًا؛ يلبس \VUgras{نظارة أنفية}.

%%K t04-tit-5 sub
\VUsaut{2.4mm}
\VUpk{10.2pt}{40}{الحلوى.}
\VUsaut{3.0mm}

%%K t04-c1-17-1 p
17. --- وخلف المدعوين، على مائدة مغطاة بـ\VUgras{مفرش}، تُبسَط
\VUgras{الحلوى} \textsuperscript{(39)}. ففي سلطانية كبيرة فاكهة:
\VUgras{خوخ} \textsuperscript{(40)} و\VUgras{كمثرى} \textsuperscript{(41)}
و\VUgras{تفاح} \textsuperscript{(42)} و\VUgras{مشمش} \textsuperscript{(43)}
و\VUgras{عنب} \textsuperscript{(44)}. وإلى جانبها \VUgras{أناناس}
\textsuperscript{(45)} و\VUgras{كعكة} \textsuperscript{(46)} جميلة
و\VUgras{زجاجات مشروب} \textsuperscript{(48)} و\VUgras{شمبانيا} \textsuperscript{(49)}،
وأكوام من \VUgras{الصحون} \textsuperscript{(47)} النظيفة كذلك.

%%K t04-c1-18-1 p
18. --- و\VUgras{ساقي الخمر} \textsuperscript{(50)} ينزع سدادة \VUgras{زجاجة}
\textsuperscript{(52)} من خمر عتيقة بـ\VUgras{مِفتاح السدادات} \textsuperscript{(51)}.
ويبدو \VUgras{الفلين} \textsuperscript{(53)} عسير النزع جدًا. ولهذا الساقي
\VUgras{منديل} \textsuperscript{(54)} على ذراعه.

%%K t04-c1-19-1 p
19. --- وبعد العشاء ينتقل الرجال إلى غرفة التدخين، وتذهب السيدات
ليستعددن للرقص.

%%K t04-tit-6 sub
\VUsaut{5.0mm}
\VUcentre{11.4pt}{0}{\textit{المشهد الثاني.}}
\VUsaut{1.6mm}
\VUpk{11.4pt}{40}{عيد ميلاد الجد.}
\VUsaut{2.6mm}

%%K t04-tit-7 sub
\VUpk{10.2pt}{40}{الزيارة.}
\VUsaut{3.0mm}

%%K t04-c2-01-1 p
1. --- مرّت عشر سنوات. شاخ والدا يوحنا، وهما يحبان أحفادهما الثلاثة حبًا
شديدًا: \VUgras{حفيدين} \textsuperscript{(2)} و\VUgras{حفيدة} \textsuperscript{(3)}.
ولأن اليوم \VUgras{عيد ميلاد الجد}، جاء الأولاد ليهنئوه.

%%K t04-c2-02-1 p
2. --- وبطرس الصغير لا يزال صغيرًا جدًا؛ عمره ثلاث سنوات فقط. يرى
\VUgras{القطة} \textsuperscript{(1)} مقبلة نحوه. يريد أن يمسح \VUgras{فروها}
الحريري الجميل و\VUgras{ذيلها} الطويل؛ ولا يخطر له أن تحت تلك
\VUgras{الكفوف} الرشيقة مخالب حادة خبيثة ربما خدشته.

%%K t04-c2-03-1 p
3. --- وأخته غابرييل، الممسكة بيده، أرزن منه بكثير، وأما مارسيل، بمعطفه
\VUgras{الكبير} \textsuperscript{(4)} و\VUgras{حزامه} \textsuperscript{(5)} الجلدي،
فيبدو رجلًا صغيرًا، على رغم السروال القصير الذي يترك \VUgras{جوربيه}
\textsuperscript{(6)} ظاهرين.

%%K t04-c2-04-1 p
4. --- وقد لبس أبوهم \VUgras{معطف} \textsuperscript{(7)} الشتاء الثقيل
بـ\VUgras{ياقته الفرائية} \textsuperscript{(8)}، وفي \VUgras{جيبه}
\textsuperscript{(9)} لِفاع. وعلى امرأته قبعتها اللبّادية المزيّنة
بـ\VUgras{ريش} \textsuperscript{(10)}، و\VUgras{سترتها} \textsuperscript{(11)}،
و\VUgras{لِفاعها الفرائي} \textsuperscript{(12)}، و\VUgras{كُمّة اليدين}
\textsuperscript{(13)}.

%%K t04-c2-05-1 p
5. --- والبرد شديد، فقد حلّ الشتاء. وظل \VUgras{الثلج} \textsuperscript{(19)}
يتساقط طوال النهار حتى ابيضّت سطوح البيوت. ومع \VUgras{الليل} يشتد القرّ
أكثر. وفي السماء يضيء \VUgras{القمر} \textsuperscript{(17)} و\VUgras{النجوم}
\textsuperscript{(18)} فيخترقان \VUgras{الظلام}.

%%K t04-tit-8 sub
\VUsaut{4.2mm}
\VUpk{10.2pt}{40}{المرض.}
\VUsaut{3.0mm}

%%K t04-c2-06-1 p
6. --- و\VUgras{الأعمى} \textsuperscript{(20)} المسكين الذي يقطع الساحة أمام
\VUgras{الصيدلية} \textsuperscript{(16)}، يقوده \VUgras{كلبه} \textsuperscript{(21)}،
تعِس جدًا. وجوستين \VUgras{الخادمة} \textsuperscript{(14)} تحمل من المطبخ
\VUgras{سلطانية} \textsuperscript{(15)} من \VUgras{شاي الأعشاب} الحار جدًا
المُحلّى تحلية سخية.

%%K t04-c2-07-1 p
7. --- و\VUgras{الجدة} \textsuperscript{(23)}، وهي تريد أن تأخذ \VUgras{نظارتها}
\textsuperscript{(26)} من على المائدة لتقرأ \VUgras{الوصفة} \textsuperscript{(27)}
التي كتبها \VUgras{الطبيب} \textsuperscript{(25)} لتوّه، تنهض من مقعدها متكئة
على \VUgras{عصاها} \textsuperscript{(24)}.

%%K t04-c2-08-1 p
8. --- وكثيرًا ما تشكو \VUgras{عللًا} صغيرة و\VUgras{دُوارًا}
و\VUgras{زُكامًا} و\VUgras{وجع أسنان}؛ ساقاها واهنتان وعيناها لم تعودا
جيدتين. ومع ذلك يطيب لها أن تمازح \VUgras{طبيبها} العجوز الذي يريد دائمًا
أن يصف لها \VUgras{شرابًا} \textsuperscript{(28)}. وهي اليوم سعيدة جدًا بأن
أسرتها الصغيرة حولها.

%%K t04-c2-09-1 p
9. --- والغرفة المحكمة الإغلاق دافئة. ففي \VUgras{المدفأة} \textsuperscript{(29)}
الرخامية تتأجج \VUgras{نار} \textsuperscript{(30)} \VUgras{بهيجة}، ويطيب المرء
نفسًا في \VUgras{ضوء} \VUgras{المصباح} \textsuperscript{(31)} الناعم الذي أُشعل
لتوّه.

%%K t04-tit-9 sub
\VUsaut{4.2mm}
\VUpk{10.2pt}{40}{الجد.}
\VUsaut{3.0mm}

%%K t04-c2-10-1 p
10. --- حتى \VUgras{الجد} \textsuperscript{(33)} نسي أنه كان مريضًا، وأن
\VUgras{الرثية} تُقعده في \VUgras{مقعده} \textsuperscript{(34)}، وأنه أصيب
بالأمس بـ\VUgras{حمى وإغماء}. لم يعد يذكر أنه أصيب الشتاء الماضي
بـ\VUgras{ذات الرئة}، وأن \VUgras{سعالًا} قاسيًا موجعًا أذاقه عذابًا شديدًا.

%%K t04-c2-10-2 p
ومع ذلك لم يبرأ إلا بعناء. وهو الآن أحسن حالًا بعض الشيء. لكن ساقه، التي
يسندها على \VUgras{وسادة} \textsuperscript{(38)} موضوعة على \VUgras{مِسند}
\textsuperscript{(39)} صغير، لا تزال تؤلمه هي أيضًا.

%%K t04-c2-11-1 p
11. --- فإذا لُفّ بعناية في \VUgras{روبه} \textsuperscript{(35)} الدافئ ذي
\VUgras{الأزرار} \textsuperscript{(36)} الصدفية الكبيرة، وكان إلى جانبه من يقرأ
له الجريدة، وهي \VUgras{اليتيمة} \textsuperscript{(40)} الصغيرة في \VUgras{قُبَّعتها}
البيضاء و\VUgras{ثوبها} \textsuperscript{(42)} الأزرق و\VUgras{عباءتها}
\textsuperscript{(41)}، فإنه لا يشكو.

%%K t04-c2-12-1 p
12. --- وفي كل سنة، حين يعيد \VUgras{التقويم} \textsuperscript{(43)}
\VUgras{تاريخ} الأول من \VUgras{ديسمبر}، ينتظر \VUgras{أحفاده} بفارغ الصبر،
فهو يعلم أنهم لن ينسوا ذلك \VUgras{التاريخ} المهم.

%%K t04-c2-13-1 p
13. --- ويذكر متأثرًا تلك الأيام البعيدة حين كان هو نفسه يذهب ليهنئ
\VUgras{المشير} الإمبراطوري المجيد، الذي يُعلَّق \VUgras{صورته}
\textsuperscript{(44)} على الجدار، وهو \VUgras{جد الأولاد الأكبر}.
\VUsaut{6.0mm}
\VUfilet{22mm}
